\chapter{86}



Joh stand auf, rückte den Stuhl neben sich zurecht und winkte Ava an den Tisch.

«Hi du, guten Morgen. Milchkaffee?»

«Sehr gerne. Danke, Joh.»

Er drehte sich luftig um und bald ertönte das Brummen der Kaffeemaschine.
Manuela hatte einen Schreibblock vor sich liegen.

«Wir haben eben den Tagesplan besprochen, Ava. Joh ersetzt mit Ciro
zusammen die Heuraufen im Stall. Dann ist dort alles bereit. Vielleicht
reicht ihnen die Zeit sogar, die alten Raufen zu entsorgen. Nur die
frischen Heuballen fehlen noch. Die werden heute Nachmittag geliefert
und eingelagert.»

Joh kam mit einem Tablett auf den Händen zurück.

«Und, Ciro, keine Sorge, wir schauen mal, was dein Bein aushält,» er
verteilte die Tassen, «und dann besprechen wir deinen Einsatz der
nächsten Tage, okay?»

«Gut, klar.»

Manuela richtete sich an Ava.

«Und was hast \emph{du} heute vor?»

«Ich steh dir zur Verfügung.»

«Schön, heute ist mein Beerentag. Ich verarbeite einen Teil unseres
Früchteberges.»

Manuela sah begeistert aus.
Louis würde mit Joh also vorwiegend alleine sein. Ihm graute. Er tanzte nur mehr wie ein Irrer zwischen zwei Welten hin und her. Ein fiktiver Ciro hatte die Hauptrolle und im Rücken lauerten Louis und die Wahrheit. 
Wie von der Decke gefallen war Valentinas Aussage zurück \emph{Vielleicht ist es gut, wenn du dich
stellst.} 

\sectionbreak

\indentedblock{Damals. Fribourg, neun.\\
Louis steht vor dem Pult der Lehrerin. Er tritt von einem Bein auf das
andere. Die Zettel auf dem Tisch wagt er nicht anzusehen. Jedenfalls
nicht direkt. Sie hat sie gesammelt. Er hat lange nachgedacht. Jedes Mal
nach ganz bösen Schimpfwörtern gesucht. Nicolas ist der dämlichste Junge
in der Klasse. Louis geht ihm aus dem Weg. Er hat Angst vor ihm.

Die Lehrerin blickt auf die Zettel.

«Also, Louis, wenn ich mir die Schrift genau ansehe, könntest du die
geschrieben haben,» sie stützt ihren Kopf auf der offenen Hand ab und
streicht sich mit den Fingern über die Wange, «hm?»

Es ist bedrohlich leer im Klassenzimmer. Ohne all die andern Kinder.

«Aber wann soll ich sie ihm denn unter sein Pult gelegt haben?»

Louis findet seine Frage schlau. Die Lehrerin nicht. Das spürt er. Ihre
Augen werden schmal. Es ist nun noch stiller, dass sich Louis ein
Geräusch herbeisehnt. Irgend etwas. Zur Not wenigstens die Pausenglocke.

Jetzt sieht er die Lehrerin einatmen.

«Wie kommst du darauf, dass Nicolas die Zettel unter seinem Pult fand?»

Sie hat recht, eigentlich kann er es nicht wissen. Zum Fundort hat sie der Klasse nichts gesagt. Louis wird sehr heiss. Der Druck 
in seinem Kopf wächst. Er will sicher nicht heulen. Aber die Tränen drücken
sich ganz von selbst in die Augen und verwischen ihm die Sicht.

«Der ist so gemein, ich,» Louis zieht die Nase hoch, «der ist viel
stärker als ich, der kann tun, was er will. Der gewinnt immer.»

«Mhm, und dann hast du ihm böse Wörter unters Pult gelegt, heimlich und
ohne deinen Namen drauf zu schreiben, ja?»

Louis nickt verlegen.

«Ja, damit er sieht, dass ihn nicht alle mögen und wie blöd er ist. Ich
kann doch nichts anderes tun bei dem.»

Die Lehrerin ist nicht böse. Wirklich nicht. Louis schaut sie zweifelnd
an.

«Das ist traurig. Dann hattest du also über die letzten Wochen eine
schwierige Zeit und wolltest mir nichts davon erzählen?»

Sie macht ein enttäuschtes Gesicht und berührt seinen Arm. Dann lächelt
sie. Und sieht wieder so aus, wie sie ihm gefällt. Die folgende Stille
hält Louis nun viel besser aus. Die Lehrerin scheint nachzudenken.

«Kannst du dir vorstellen, Louis, dass du zu mir kommst, wenn du das
nächste Mal an Nicolas gerätst?»

«Ja, vielleicht. Aber dann bin ich ein Petzer und wenn es dann noch
schlimmer wird?»

«Das kann es nicht werden, Louis, dafür sorge ich.»

Ihre Stimme klingt heller als sonst. Louis glaubt ihr. Jetzt fällt es
ihm leicht, sich zu entschuldigen. Sie nickt anerkennend.

Und als sie dann vor der Klasse steht. Gleich am nächsten Morgen,
schafft er es, sie anzusehen. So, als habe er mit den bösen Nachrichten
nichts zu tun.

Sie macht es kurz. Nicolas sieht zufrieden aus. Mit ihm hat sie bestimmt
schon gesprochen.

«Nun, Kinder, das Rätsel ist gelöst. Ich weiss nun, wer Nicolas die
hässlichen Zettel unters Pult gelegt hat. Das Kind hat es zugegeben, mir
den Grund erklärt und sich entschuldigt. Damit ist die Sache erledigt.»

Louis schielt verstohlen zu Nicolas hinüber. Der nickt.

Dann erklärt die Lehrerin, was «anonym» bedeutet. Daraus entwickelt sich
ein interessantes Gespräch. Louis liebt Klassengespräche. Doch diesmal
hört er still zu und meldet sich nicht. Die Lehrerin sieht so schön aus.}

\sectionbreak

Die weiteren Gespräche blendete Louis aus. Immer wieder nahm er Avas Blicke wahr. Die nächtliche
Szene war ihm peinlich. Dann fiel ihm Valentina ein. Er
konnte sie erst später anrufen. Was sie ihm wohl Dringendes sagen
wollte? Doch vielleicht war es klug, jetzt gleich mit ihr zu sprechen.
Vielleicht hatte sie wichtige, neue Informationen. Louis’ Gedanken
erschienen ihm wie Gegner, die ihren Raum verteidigten. Er hätte schreien mögen. Je mehr er sich anstrengte, 
desto weniger vertrugen sie sich. Es reichte. Er stand abrupt auf und stellte seinen Stuhl zurecht.

Alle drei schauten ihn verwundert an.

«Ich habe noch kurz einen Anruf zu erledigen. Ist es okay, wenn ich in
einer Viertelstunde zu dir in den Stall komme, Joh.»

«Ja, passt,» Joh hob den Daumen, «und zieh dir feste Schuhe an. Wir
brauchen Halt.»

Jetzt grinste er.

\sectionbreak

Louis humpelte die Treppe hoch. Er schloss die Tür des
\emph{Schäferlofts} und blieb im Raum stehen. Wann genau
hatte ihm Valentina eigentlich die Nachricht geschickt? Er suchte nach
der Zeitangabe auf dem Handy. Es war vor dem Frühstück gewesen. Er hatte
gerade geduscht. Das Gerät hatte irgendwo bei seinen Sachen gelegen. Es
war zwar auf lautlos eingestellt. Trotzdem, er wurde zu nachlässig.

Louis setzte sich aufrecht auf sein Bett und wählte Valentinas Nummer.
Ohne Begrüssung erreichte ihn ein Schwall von Worten. Sie war
ungehalten, sprach laut und ihre Stimme verriet Ärger.

«Du hast mich gestern einfach plötzlich abserviert. Was läuft denn da
bei dir? Ich mache mir echt Sorgen um dich, Louis, warum gehst du nicht
zur Polizei, erzählst, was du gesehen hast. Und klar, dein Verschwinden
war blöd, eine Kurzschlussreaktion eben, die werden das bestimmt - »

«Du hast völlig recht, Valentina. Es tut mir leid. Ich hatte keine
andere Wahl, als den Anruf zu beenden. Und ja, ich denke auch, ich muss - ich, ich bin dran, das kannst du mir glauben. Wolltest du deshalb
mit mir reden?»

«Auch,» ihre Stimme bekam einen feineren Klang, «ach,
Louis, weisst du, ich, das macht mich alles so traurig, hm. Was ich dir
eigentlich sagen wollte, ist was anderes. Da hat sich etwas geklärt. Der gesuchte
Mann, der in der Nacht auf der Polizei auftauchte, wurde gefunden. Ob er
mit Giorgias Unfall etwas zu tun hat, wird offenbar untersucht - »

Louis unterbrach sie erneut. Ihm blieb wenig Zeit.

«Danke, Valentina, ja, das habe ich gelesen.»

«Gut, dann weisst du das bereits. Aber nun das Wichtigste, Louis, ich
habe mich gerade bei den Salas nach Giorgia erkundigt. Sie macht zwar
Fortschritte. Die OP war erfolgreich, nur,» Louis wollte nicht hören,
was folgte, «sie spürt ihre Beine nicht, also, noch nicht, das könne
wieder kommen, sagen die Ärzte, aber sicher, so ganz sicher ist es
nicht,» Valentina rang mit sich, «ihre Mutter hat wieder geweint. Es hat
mir das Herz zerrissen. Giorgia darf ab morgen Besuch bekommen. Ich
werde mit Maria und Franco hingehen. Der Arzt meint, sie braucht nun
vertraute Menschen um sich, es werde ihr bei der Genesung helfen. So
richtig wach wird sie nicht sein, sagt er. Sie döse mehrheitlich vor
sich hin. Aber auch in diesem Zustand seien vertraute Stimmen für sie
wichtig und heilsam.»

Louis stürzte in sich zusammen. Er brachte kein Wort heraus. Für einen
Moment glaubte er, sein Herz gebe auf.

«Louis, Louis? - bist du noch da?»


